\documentclass[11pt,a4paper,oneside,openany]{book}
\usepackage[a4paper,left=34mm,right=34mm,top=28mm,bottom=30mm]{geometry}
\usepackage[T1]{fontenc}
\usepackage{newtxtext,newtxmath}
\usepackage{microtype}
\usepackage{amsmath,amssymb}
\usepackage{booktabs,longtable,tabularx,array}
\usepackage[dvipsnames]{xcolor}
\usepackage{listings}
\usepackage{enumitem}
\usepackage{fancyhdr}
\usepackage{titlesec}
\usepackage{hyperref}
\usepackage{bookmark}
\usepackage{parskip}
\usepackage{graphicx}
\usepackage{caption}
\usepackage{etoolbox}

\definecolor{PySigRed}{HTML}{7A1731}
\definecolor{PySigInk}{HTML}{252525}
\definecolor{CodeBg}{HTML}{F5F2EE}
\definecolor{CodeRule}{HTML}{C9BFB6}
\hypersetup{colorlinks=true,linkcolor=PySigRed,urlcolor=PySigRed,citecolor=PySigRed,pdftitle={PySig 0.4.3 User Manual},pdfauthor={PySig Project}}
\setlength{\parindent}{0pt}
\setlength{\parskip}{5pt}
\setcounter{secnumdepth}{3}
\setcounter{tocdepth}{2}
\titleformat{\chapter}[display]{\bfseries\color{PySigRed}\huge}{\chaptername\ \thechapter}{6pt}{\Huge}
\titleformat{\section}{\bfseries\color{PySigRed}\Large}{\thesection}{0.7em}{}
\titleformat{\subsection}{\bfseries\color{PySigInk}\large}{\thesubsection}{0.7em}{}
\pagestyle{fancy}\fancyhf{}\fancyhead[L]{\small PySig 0.4.3 User Manual}\fancyhead[R]{\small Intersection and Network Modelling}\fancyfoot[C]{\thepage}\renewcommand{\headrulewidth}{0.4pt}
\lstdefinelanguage{yaml}{keywords={true,false,null},sensitive=false,comment=[l]{\#},morestring=[b]",morestring=[b]'}
\lstdefinestyle{pysig}{basicstyle=\ttfamily\small,backgroundcolor=\color{CodeBg},frame=single,rulecolor=\color{CodeRule},breaklines=true,showstringspaces=false,columns=fullflexible,xleftmargin=4pt,xrightmargin=4pt,aboveskip=8pt,belowskip=8pt}
\lstset{style=pysig}
\newcommand{\code}[1]{\texttt{#1}}
\newcommand{\PySig}{\textbf{PySig}}
\newenvironment{modelnote}{\begin{quote}\small\itshape}{\end{quote}}
\begin{document}
\frontmatter
\begin{titlepage}\centering\vspace*{35mm}{\color{PySigRed}\bfseries\Huge PySig 0.4.3\par}\vspace{8mm}{\bfseries\LARGE User Manual\par}\vspace{7mm}{\large Constructing, assigning and optimising intersection and network models\par}\vfill{\itshape Australian left-hand traffic conventions\par}YAML configuration · lane-based queues · fixed OD demand · signal optimisation\par\vspace{10mm}\today\end{titlepage}
\chapter*{Preface}\addcontentsline{toc}{chapter}{Preface}
This manual is the source-controlled documentation for \PySig{} version 0.4.5. It is written as a practical model-building guide rather than only a schema reference. The central workflow is: draw the physical intersection, inventory lanes, define lane-to-lane movements, define signal control, convert observed turning flows into fixed OD demand, validate the structure, run a baseline, then optimise.

\begin{modelnote}
Fixed OD demand is an invariant. Assignment may redistribute an OD total among permitted routes and lanes, but it must not create, suppress or move demand between zone pairs.
\end{modelnote}

PySig is a developing analytical package. Results should be calibrated and validated before use in formal operational-modelling, design, approval or policy decisions.

\tableofcontents
\mainmatter
\chapter{Getting started}
\section{Requirements}
PySig is intended to run in a restricted Python environment. The core project requires Python 3.10 or later and PyYAML. The test suite uses pytest.
\begin{lstlisting}[language=bash]
conda install pyyaml pytest
pytest -v
python main.py
\end{lstlisting}

\section{Recommended repository layout}
\begin{lstlisting}
PySig/
|-- main.py
|-- traffic_model/
|-- configs/
|-- tests/
|-- docs/
|   `-- manual/
|       |-- pysig_manual.tex
|       |-- chapters/
|       `-- examples/
`-- .github/workflows/
    `-- build-manual.yml
\end{lstlisting}

\section{Run modes}
For an isolated optimised intersection, use:
\begin{lstlisting}[language=yaml]
run:
  mode: optimise_intersection
\end{lstlisting}
For a network assignment without signal optimisation, use \code{network}. For coordinated optimisation use \code{optimise_network}.

\chapter{The PySig model in one page}
Two chains are fundamental.

\section{Physical and control chain}
\begin{center}
\textit{intersection $\rightarrow$ lane $\rightarrow$ movement $\rightarrow$ signal group $\rightarrow$ phase $\rightarrow$ controller}
\end{center}
A lane is physical storage. A movement is a permitted lane-to-lane path. A signal group controls one or more movements. A phase activates signal groups. A controller repeats the phase sequence.

\section{Demand chain}
\begin{center}
\textit{OD demand $\rightarrow$ route $\rightarrow$ movement(s) $\rightarrow$ lane queue(s)}
\end{center}
The OD matrix fixes how much traffic travels between each pair of zones. Assignment decides which permitted routes carry that fixed demand.

\section{Shared-lane kernel}
Movements using the same \code{from_lane_id} share one physical queue and storage constraint. Therefore
\begin{equation}
q_{\mathrm{lane}}=\sum_{m\in M_{\mathrm{lane}}}q_m.
\end{equation}
Lane saturation flow is applied once and shared by eligible movements. This prevents two movements on one lane independently consuming the lane's full physical capacity.

\chapter{Build an intersection from scratch}
\section{Step 1: mark the arms on a lane diagram}
Begin with the physical layout, not the YAML. Assign a consistent arm identifier to each approach. A four-arm intersection might use \code{A1}, \code{A2}, \code{A3}, and \code{A4}. Keep the convention stable across lanes, movements and working notes.

\section{Step 2: inventory every physical lane}
Create a worksheet before coding.
\begin{center}\begin{tabularx}{\textwidth}{llllX}\toprule
Lane ID & Arm & Type & Length & Notes\\\midrule
A1\_1 & A1 & regular & 220 m & through/left shared\\
A1\_2 & A1 & regular & 120 m & right turn\\
A2\_3 & A2 & regular & 220 m & parent lane\\
A2\_3L & A2 & short & 80 m & short lane from A2\_3\\\bottomrule
\end{tabularx}\end{center}

\section{Step 3: create the intersection envelope}
\begin{lstlisting}[language=yaml]
intersections:
  - id: I1
    controller_id: C1
    movement_ids:
      - T_WE
      - L_WN
      - R_WS
\end{lstlisting}
Explicit \code{intersection_id} values are recommended even for a single intersection because the model can later be incorporated into a network without renaming objects.

\section{Step 4: build lanes before movements}
Do not use a drawing label such as \code{1} in \code{from_lane_id} if the YAML lane is called \code{A1_1}. References must match the complete object ID exactly.

\chapter{Lanes, short lanes and movements}
\section{Regular lanes}
\begin{lstlisting}[language=yaml]
lanes:
  - id: A1_1
    intersection_id: I1
    physical_length_m: 220
    saturation_flow:
      method: direct
      value_pcu_per_hour: 1900
\end{lstlisting}
If storage is derived from physical length, the basic relationship is
\begin{equation}
K_{\mathrm{lane}}=\frac{L_{\mathrm{lane}}}{l_{\mathrm{PCU}}}.
\end{equation}
With a 7.35 m PCU, a 220 m lane stores approximately 29.9 PCU.

\section{Short lanes}
A short lane references an existing parent lane. Multiple short lanes may share the same parent.
\begin{lstlisting}[language=yaml]
  - id: A2_3L
    intersection_id: I1
    parent_lane_id: A2_3
    physical_length_m: 80
    saturation_flow:
      method: rr67_unopposed
      lane_width_m: 3.2
\end{lstlisting}

\section{Movements are lane-to-lane paths}
A movement should be read as
\begin{equation}
\textit{movement}=\textit{from lane}\rightarrow\textit{to lane}.
\end{equation}
\begin{lstlisting}[language=yaml]
movements:
  - id: T_WE
    intersection_id: I1
    from_lane_id: A1_1
    to_lane_id: A3_4
    movement_type: through
    signal_group_id: SG_A1
    free_flow_time_s: 8
\end{lstlisting}
\code{to_lane_id} is optional for a movement that exits the model, but when supplied it should reference a real lane.

\section{Several movements on one lane}
If \code{T_WE} is 500 PCU/h and \code{L_WN} is 100 PCU/h and both use \code{A1_1}, the physical lane demand is 600 PCU/h. PySig retains movement composition but calculates one physical queue, delay process and storage constraint for \code{A1_1}.

\chapter{Signal control and effective service}
\label{chap:signal-control}

\section{Controllers, phases and groups}
A controller defines cycle time, offset and ordered phases.  Each phase contains one or
more active signal groups.  Movements reference signal groups rather than phases
directly.  This permits one signal group to be active in more than one phase.

\section{Displayed green and intergreen}
For signal group $j$, displayed signal green is the duration for which the group is
active.  Intergreen belongs to transition/clearance time and is not automatically
counted as signal green.

For a simple non-overlapping phase sequence,
\begin{equation}
C=\sum_p(G_p+I_p),
\end{equation}
where $C$ is cycle time, $G_p$ is displayed phase green and $I_p$ is the intergreen
after phase $p$.

\section{Continuous signal-group runs and phase overlap}
If a signal group is active in consecutive phases, PySig treats the group as
continuously active across the internal phase boundary.  For example:

\begin{verbatim}
Phase P_A: SG_A, SG_B
Phase P_B:       SG_B, SG_C
\end{verbatim}

SG_B has one continuous signal run.  The intergreen after P_A is not inserted into
SG_B's service simply because the controller changes phase.  Ending displacement and
clearance treatment apply when the continuous SG_B run actually ends.

\section{Start displacement and end displacement}
PySig distinguishes displayed signal green $G$ from engineering effective green $g$.
With start displacement $d_s$ and end displacement $d_e$,
\begin{equation}
g=G-d_s+d_e,
\end{equation}
bounded to the interval $0\leq g\leq C$.

For 55 s of displayed green, 3 s start displacement and 3 s end displacement,
\begin{equation}
g=55-3+3=55\text{ s}.
\end{equation}
The duration is unchanged, but the effective service interval is displaced in time.
An asymmetric 3 s start displacement and 2 s end displacement gives 54 s of effective
green.

\section{Capacity implication}
The nominal fixed-time capacity is
\begin{equation}
c_n=s\frac{g}{C},
\end{equation}
where $s$ is lane saturation flow.  The corresponding nominal degree of saturation is
\begin{equation}
x_n=\frac{q}{c_n}=\frac{qC}{sg}.
\end{equation}
These relationships are consistent with conventional signal-capacity definitions.
Operational accessible capacity may be lower after lane interactions are considered.

\section{Priority and yielding movements}
A yielding movement may identify one or more priority movements and an active capacity
factor.  When no priority movement is active, the factor is 1.0.  While a configured
priority stream is effectively active, available yielding service is reduced by the
configured factor.  The factor is applied within the shared physical-lane supply; it
does not create separate full lane capacity for each movement.

The current factor formulation is an explicit planning approximation.  It is not a
full gap-acceptance model involving critical gap, follow-up headway, pedestrian
conflicts or opposing platoon structure.

\chapter{Zones, routes and fixed OD demand}
\section{Convert turning counts to OD demand}
For an isolated intersection, approach and exit zones are a transparent starting point. If the west approach contains 500 PCU/h through, 80 PCU/h left and 70 PCU/h right, define three OD movements rather than treating the three turns as alternative routes to the same destination.
\begin{lstlisting}[language=yaml]
routes:
  - id: R_W_E
    origin_zone_id: Z_W
    destination_zone_id: Z_E
    movement_ids: [T_WE]
  - id: R_W_N
    origin_zone_id: Z_W
    destination_zone_id: Z_N
    movement_ids: [L_WN]
  - id: R_W_S
    origin_zone_id: Z_W
    destination_zone_id: Z_S
    movement_ids: [R_WS]

origin_destination_matrix:
  values:
    Z_W:
      Z_E: 500
      Z_N: 80
      Z_S: 70
\end{lstlisting}
Thus the west approach total is $500+80+70=650$ PCU/h.

\section{OD conservation}
For every OD pair $ij$,
\begin{equation}
\sum_{r\in R_{ij}} q_r=Q_{ij}.
\end{equation}
The matrix total is fixed. Assignment changes only the route distribution.

\section{Assignment convergence}
PySig uses a logit target allocation and method of successive averages. A useful representation is
\begin{equation}
w_r=\exp[-\theta(c_r-c_{\min})]
\end{equation}
and
\begin{equation}
q_r^n=q_r^{n-1}+\frac{1}{n}(\hat q_r^n-q_r^{n-1}).
\end{equation}
Convergence is assessed using maximum route-flow change. OD conservation is an accounting identity, not a behavioural convergence criterion.

\chapter{Saturation flow and lane capacity}
\label{chap:saturation-flow}

\section{Definition}
Saturation flow is the maximum queue-discharge flow rate at a stop line under continuous
effective service, expressed in PCU/h of effective green.  It is not the same as lane
capacity over a complete signal cycle.  Nominal capacity is obtained by applying the
effective-green ratio.

\section{Direct input}
A measured or otherwise justified saturation flow may be entered directly:
\begin{verbatim}
saturation_flow:
  method: direct
  value_pcu_per_hour: 1900
\end{verbatim}
Field measurement is preferable where suitable observations are available because
saturation flow strongly affects capacity, DoS, queue and delay.

\section{RR67-style estimation}
PySig also implements an RR67-style unopposed saturation-flow estimate.  Supported
inputs include lane width, gradient, nearside-lane status, turning proportion,
turning radius and a local calibration factor.  The software output should expose both
the final saturation-flow value and the inputs used to obtain it.

Example:
\begin{verbatim}
saturation_flow:
  method: rr67_unopposed
  lane_width_m: 3.2
  gradient_percent: 1.5
  is_nearside_lane: false
  turning_radius_m: 18.0
  local_factor: 1.0
\end{verbatim}

\section{Dynamic turning proportion}
If neither \code{turning\_proportion} nor \code{turning\_proportion\_override} is supplied,
PySig derives the current lane turning proportion from assignment.  This avoids treating
a behavioural lane-flow property as permanently fixed before assignment.

A historical \code{turning\_proportion} remains a fixed-input override for compatibility.
The preferred explicit form for a calibrated fixed assumption is
\code{turning\_proportion\_override}.

\section{Turning radius}
Turning radius is supplied in metres.  Where one turn path dominates the lane,
\code{turning\_radius\_m} may be stored in the lane saturation-flow block.  Where several
turning movements provide movement-specific radii, PySig can derive a turning-flow-
weighted radius.  This is an assignment-responsive approximation rather than a claim
that the weighted radius has independent geometric meaning.

\section{Base and accessible capacity}
RR67 returns a base saturation-flow estimate $s$.  Effective green gives nominal
capacity
\begin{equation}
c_n=s\frac{g}{C}.
\end{equation}
Short lanes, shared-lane blocking and priority effects can then reduce or redistribute
the service accessible to individual lanes.  For auditability, these interactions do
not overwrite the base RR67 value.

\section{Calibration}
A calculated saturation flow is a prediction.  Site-specific measured discharge should
be used to calibrate or override prediction where adequate data exist.  Queue discharge,
start loss/end gain and saturation flow are observable field quantities and should be
checked together for important models.

\chapter{Signal optimisation}
\label{chap:optimisation}

\section{Nested problem structure}
Optimisation and traffic assignment solve different problems.  Assignment redistributes
fixed demand for one fixed signal plan.  Optimisation compares alternative signal
plans, each of which requires its own traffic assignment.

\begin{verbatim}
signal candidate
  -> route/lane assignment
  -> dynamic capacity update
  -> lane performance
  -> feasibility and objective
  -> next candidate
\end{verbatim}

\section{Sequential search}
For sequential optimisation, the supplied controller must be the centre of the search.
Candidate zero is the current controller plan:
\begin{verbatim}
cycle  = current cycle
splits = current phase durations
offset = current offset
\end{verbatim}
The current plan is therefore always evaluated before alternatives.  Nearby candidates
are created by transferring the configured split increment between phases while
preserving cycle constraints and minimum greens.  This is preferable to taking the
first combinations from a large arbitrary enumeration because realistic multi-phase
controllers can have many thousands of combinatorial split possibilities.

\section{Minimum green and cycle reconciliation}
A candidate is timing-valid only if every phase respects the configured minimum green
and the controller cycle reconciles with phase greens and intergreens.  For a simple
non-overlap cycle,
\begin{equation}
\sum_p G_p + \sum_p I_p = C.
\end{equation}
The optimiser should not generate candidates that violate its own timing constraints.

\section{Candidate evaluation}
For each candidate, PySig performs route assignment and then evaluates the configured
objective, normally network/intersection delay.  Candidate diagnostics should identify
cycle, phase splits, offset, delay, convergence status and feasibility.

PySig should distinguish hard invalidity (for example, an impossible signal plan,
numerical failure or conservation failure) from poor operational performance.  High
DoS, residual queue and short-lane storage usage are performance measures that should
remain visible to the optimisation objective rather than being silently hidden.

\section{Search limits}
\code{max\_candidates\_per\_controller} is a computational limit.  When the candidate
space is larger than this limit, candidate ordering becomes material.  Starting from
the current plan and exploring nearby candidates gives this limit an operationally
meaningful interpretation.

\section{Convergence}
The optimisation outer loop should consider stability of both signal plans and assigned
flows.  With assignment-responsive RR67 and nested short-lane capacity, changes in
assigned flow can also alter lane capacity.  Capacity-change metrics are therefore
useful additional diagnostics even where the formal stopping criterion remains based on
flow change and signal-plan stability.

\chapter{Multi-intersection networks}
A network links the exit movement of one intersection to the approach movement of another. Link free-flow travel time is
\begin{equation}
t_{\mathrm{link}}=\frac{3.6L}{v},
\end{equation}
where $L$ is metres and $v$ is km/h.
\begin{lstlisting}[language=yaml]
network_links:
  - id: LINK_MAIN
    from_intersection_id: I1
    to_intersection_id: I2
    upstream_movement_id: M1_MAIN
    downstream_movement_id: M2_MAIN
    length_m: 500
    speed_kmh: 60
    storage_pcu: 65
    dispersion: 0.10
\end{lstlisting}
At 500 m and 60 km/h, free-flow travel time is 30 seconds. Controllers use a shared network clock, so offsets become meaningful in a connected network.

\section{Sequential network optimisation}
The optimiser varies one enabled controller while holding the others fixed, reruns assignment and network simulation for every candidate, retains the best feasible plan, then proceeds to the next controller. Outer iterations repeat until signal plans and route flows stabilise or the configured limit is reached.

\chapter{Results, interpretation and model testing}
\label{chap:results-testing}

\section{Result hierarchy}
Results should be interpreted from the lane upward.  Movement results explain movement
composition and service; lane results represent the physical queue/capacity unit;
intersection and network summaries aggregate those results.

\section{Core lane outputs}
Important quantities include:
\begin{itemize}
\item demand in PCU/h;
\item base saturation flow in PCU/h of effective green;
\item displayed and effective green;
\item nominal capacity and nominal DoS;
\item accessible/operational capacity and operational DoS;
\item arrivals and departures;
\item maximum and residual queue;
\item queue estimate in metres;
\item total queue-area delay and mean queueing delay;
\item finite-storage utilisation and overflow;
\item LOS together with its basis/status;
\item capacity method and audit inputs.
\end{itemize}

\section{Nominal versus operational degree of saturation}
Nominal DoS is
\begin{equation}
x_n=\frac{q}{c_n}.
\end{equation}
Operational DoS instead uses accessible capacity $c_a$:
\begin{equation}
x_a=\frac{q}{c_a}.
\end{equation}
The two values may differ materially for short lanes, shared lanes, yielding movements
or other constrained lane systems.  Reports should not label one generic value simply
``degree of saturation'' without identifying the denominator.

\section{Residual queue, unserved demand and overflow}
These outputs are different:
\begin{itemize}
\item \textbf{residual/final queue}: admitted traffic still physically queued at the
      analysis end;
\item \textbf{unserved demand}: admitted arrivals not departed over the reporting
      period;
\item \textbf{storage overflow}: demand prevented from entering finite lane storage.
\end{itemize}
A zero overflow does not imply zero residual queue.

\section{Output precision}
Internal calculations retain floating-point precision.  Human-readable operational
outputs may be rounded to two decimal places, whereas conservation and convergence
diagnostics should retain higher precision.  Rounding belongs in serialisation, not in
the assignment, queue or optimisation calculations.

\section{Regression testing}
PySig's test suite should protect both software behaviour and engineering invariants.
Core invariants include:
\begin{itemize}
\item OD flow conservation;
\item lane queue equals movement subqueue composition where applicable;
\item lane delay equals allocated movement delay;
\item physical lane saturation flow is not duplicated across movements;
\item displayed green excludes intergreen;
\item effective green honours displacement and overlap;
\item dynamic turning proportion matches assigned flows;
\item short-lane hierarchy contains no cycles;
\item child lane length is shorter than parent lane length;
\item nested tree lane-capacity allocations sum to tree capacity.
\end{itemize}

\section{Validation hierarchy}
Validation should be undertaken at four levels:
\begin{enumerate}
\item unit and equation tests against hand calculations;
\item conservation and structural regression tests;
\item comparison between simulation and analytical methods for compatible cases;
\item field or independent-software validation for consequential applications.
\end{enumerate}

\chapter{Troubleshooting}
\begin{longtable}{p{0.29\textwidth}p{0.62\textwidth}}\toprule
\textbf{Symptom} & \textbf{Check}\\\midrule\endhead
Missing lane reference & \code{from_lane_id}, \code{to_lane_id}, and \code{parent_lane_id} must match complete lane IDs exactly.\\
Cycle mismatch & Sum every phase duration and intergreen and compare with \code{cycle_time_s}.\\
OD pair has no route & Add a valid route for every positive matrix cell or remove the demand.\\
Assignment does not converge & Review route alternatives, $\theta$, signal timings and the iteration limit; do not simply accept unconverged results.\\
Unexpected low delay & Check that all real turning movements and OD demand have been represented, and inspect unserved demand and storage overflow.\\
Old tests fail after upgrade & Extract the release into a clean directory rather than overlaying files from earlier schema versions.\\\bottomrule
\end{longtable}

\appendix
\chapter{YAML reference}
\begin{longtable}{p{0.35\textwidth}p{0.55\textwidth}}\toprule
\textbf{Field} & \textbf{Purpose}\\\midrule\endhead
\code{run.mode} & \code{network}, \code{optimise_intersection}, or \code{optimise_network}.\\
\code{settings.pcu_length_m} & Physical metres represented by one PCU; 7.35 m is the established project default.\\
\code{lane.parent_lane_id} & Parent regular lane for a short lane.\\
\code{movement.from_lane_id} & Physical lane whose queue contains the movement.\\
\code{movement.to_lane_id} & Immediate receiving lane; optional for model exits.\\
\code{controller.offset_s} & Controller offset on the shared network clock.\\
\code{assignment.theta} & Logit route-cost sensitivity.\\
\code{network.time_slice_s} & Time resolution of the network/lane simulation.\\
\code{unconverged_candidate_policy} & Use \code{reject} for formal optimisation.\\\bottomrule
\end{longtable}

\chapter{Equation summary}
\label{app:equations}

\begin{align}
\text{Effective green:}\qquad
 g &= G-d_s+d_e,\\[4pt]
\text{Nominal capacity:}\qquad
 c_n &= s\frac{g}{C},\\[4pt]
\text{Nominal DoS:}\qquad
 x_n &= \frac{q}{c_n}=\frac{qC}{sg},\\[4pt]
\text{Operational DoS:}\qquad
 x_a &= \frac{q}{c_a},\\[4pt]
\text{Fluid queue:}\qquad
 \frac{dQ}{dt} &= a-\mu,\\[4pt]
\text{Queue progression:}\qquad
 Q_1 &= Q_0+(a-\mu)\tau,\\[4pt]
\text{Clearance time:}\qquad
 t_c &= \frac{Q_0}{\mu-a}\quad(\mu>a),\\[4pt]
\text{Queue-area delay:}\qquad
 D &= \frac{Q_0+Q_1}{2}\tau,\\[4pt]
\text{OD conservation:}\qquad
 \sum_{r\in R_w} f_r &= q_w,\\[4pt]
\text{Logit share:}\qquad
 P_r &= \frac{e^{-\theta C_r}}{\sum_k e^{-\theta C_k}},\\[4pt]
\text{MSA update:}\qquad
 f_r^{(n)} &= f_r^{(n-1)}+\frac{1}{n}\left(f_r^*-f_r^{(n-1)}\right),\\[4pt]
\text{Assigned turning share:}\qquad
 p_t &= \frac{\sum_{m\in T_l}q_m}{\sum_{m\in M_l}q_m},\\[4pt]
\text{Lane storage:}\qquad
 K_i &= \frac{L_i}{l_{PCU}},\\[4pt]
\text{Short-lane cycle boost:}\qquad
 \Delta c_i &= \min\left(c_{n,i},K_i\frac{3600}{C}\right),\\[4pt]
\text{Lane-tree capacity:}\qquad
 c_{tree} &= \min\left(\sum_i c_{n,i}, c_{n,root}+\sum_{i\ne root}\Delta c_i\right),\\[4pt]
\text{Shared-lane pre-block expectation:}\qquad
 E[N] &= \frac{p_g}{p_b}\quad(p_b>0),\\[4pt]
\text{Network travel time:}\qquad
 t_f &= \frac{3.6L}{v}.
\end{align}

\chapter{First-run acceptance checklist}
\begin{itemize}[leftmargin=*]
\item Every lane ID is unique.
\item Every parent lane exists.
\item Every movement's from lane exists.
\item Every supplied to lane exists.
\item Every movement belongs to the intended intersection.
\item Every signal group appears in the intended phase sequence.
\item Phase durations plus intergreens equal the initial controller cycle.
\item Every positive OD pair has at least one valid route.
\item Turning counts have been converted to the correct OD destinations.
\item Saturation-flow sources and local factors are documented.
\item OD conservation passes.
\item Assignment converges for the selected plan.
\item Lane and link storage constraints pass.
\item The complete pytest suite passes.
\end{itemize}

\backmatter
\chapter*{Source framework}\addcontentsline{toc}{chapter}{Source framework}
PySig terminology and modelling structure have been informed by LinSig lane, movement, signal-group and phase concepts; Western Australian operational-modelling practice; and RR67-style saturation-flow calculation. These sources do not certify PySig or imply feature-for-feature equivalence.
\end{document}
